The 4096-molecule liquid agreement was recounted with the Nguyen and
Molinero 2015 rule as the default.
That rule gives an ice or clathrate label only when a molecule has
exactly four neighbours inside 3.5~\AA.
It averages $q_3$ over those four nearest neighbours.
It also applies the interfacial neighbour test in Table~1 of that paper.
We read the shell count from Table~1 of that paper.
A molecule with fewer than four neighbours inside the cutoff is liquid.
A molecule with more than four neighbours inside the cutoff is liquid.
The published rule stops at that shell test.
The four nearest neighbours are then the only bonds that enter $q_3$.
Each of those four bonds is staggered or eclipsed or neither.
A staggered bond has a correlation of at most $-0.8$.
An eclipsed bond has a correlation from $-0.35$ to $0.25$, with both
ends included.
The interfacial row of Table~1 looks at the neighbours after the bond
counts.
That row does not relabel a molecule whose shell length is not four.
The ovito-chill+ option keeps a shell of four neighbours inside 3.5~\AA.
Under that option the average for $q_3$ runs over every neighbour in the
shell.
Under that option the neighbour test is off.
Under that option a correlation equal to $-0.8$ is not staggered.
Under that option a correlation equal to $-0.35$ is not eclipsed.
Under that option a correlation equal to $0.25$ is not eclipsed.
Plain CHILL is a separate classifier.
Its thresholds are the ones already in the code.
The recount uses the first frame of the TIP4P dump.
The selected atoms have LAMMPS type 2.
That frame contains 250 such molecules.
The cutoff in the recount is 3.5~\AA.
The earlier shell count was a bug.
It scored the four nearest neighbours after the bond list had been cut
to four, even when a fifth neighbour lay inside the cutoff.
That earlier rule reported 12 interclathrate molecules.
It also reported 238 liquid molecules.
The published rule reports 9 interclathrate molecules.
It also reports 241 liquid molecules.
Three molecules change class, and each of them moves from interclathrate
to liquid.
Their LAMMPS ids are 85 and 664.
The remaining id is 679.
Molecule 85 is index 28 in the type 2 list.
Molecule 664 is index 221.
Molecule 679 is index 226.
Each of those molecules has five neighbours inside 3.5~\AA.
The other interclathrate molecules from the earlier rule have four
neighbours inside the cutoff.
Those molecules stay interclathrate.
No molecule leaves the liquid class.
The arithmetic from the earlier interclathrate count is 12 minus 3,
which is 9.
The liquid count is 238 plus 3, which is 241.
For molecule 85 the four nearest correlations include 0.582, which lies
above the eclipsed window.
The other three correlations of molecule 85 lie inside that window.
For molecule 664 the four nearest correlations include 0.317, which lies
above the eclipsed window.
The other three correlations of molecule 664 lie inside that window.
For molecule 679 the four nearest correlations include $-0.788$, which
lies between the staggered bound and the eclipsed window.
The other three correlations of molecule 679 lie inside the eclipsed
window.
None of the three molecules has a staggered bond among those four.
The neighbour test looks for staggered bonds, so it does not fire for
them.
Their earlier interclathrate label came from three eclipsed bonds.
The published rule writes liquid as soon as the shell length is not
four.
The fifth neighbour is enough to end the ice decision.
The ovito-chill+ option keeps the requirement of four neighbours inside
3.5~\AA.
Under that option, $q_3$ is averaged over every neighbour inside the
cutoff, and the neighbour test is left off.

The 4096-molecule liquid agreement in the macros is that recount under
the 2015 default.
The ovito-chill+ option is reported apart from those macros.
